\chapter{Module 2 : vivier de CV et gestion des offres d'emploi}

\section*{Introduction}
\addcontentsline{toc}{section}{Introduction}
Ce deuxième module matérialise le premier cœur opérationnel du produit. Après avoir sécurisé l'accès et structuré les rôles, il fallait construire la base documentaire et métier sur laquelle reposera toute la suite de la plateforme. Dans \projectname{}, cela signifie deux choses : constituer un vivier de CV exploitable et formaliser les besoins de recrutement à travers des offres structurées. Sans ce module, il n'y aurait ni pipeline candidat fiable, ni matching pertinent, ni recherche sémantique, ni pipeline RAG réellement utile.

Ce chapitre marque le passage d'un système simplement authentifié à un système capable de manipuler une matière métier réelle. Les CV sont hétérogènes, souvent peu structurés et parfois redondants. Les offres doivent quant à elles rester lisibles pour les recruteurs tout en devenant suffisamment normalisées pour servir aux futurs traitements IA.

\section{Sprint 2 : objectif et périmètre}

\subsection{Objectif du sprint}
L'objectif de ce sprint est de transformer deux objets bruts en objets intelligibles par le système :
\begin{itemize}
  \item le CV, qui arrive sous forme de document PDF ou DOCX ;
  \item l'offre d'emploi, qui doit être exprimée de manière suffisamment structurée pour être comprise à la fois par les recruteurs et par les services IA.
\end{itemize}

Le sprint met donc en place une page de gestion du vivier, l'import de CV avec extraction automatique, l'indexation sémantique, la détection des doublons, la recherche et l'export du vivier, la création et la gestion des offres ainsi que la préparation du matching et du futur pipeline candidat.

Ce module est stratégique parce qu'il constitue la jonction entre l'interface et l'intelligence. C'est à partir d'ici que les données commencent à devenir une matière première exploitable pour le matching, le RAG et l'agent conversationnel.

\subsection{Backlog du sprint}
\begin{table}[H]
  \centering
  \caption{Backlog fonctionnel du module vivier de CV et offres}
  \begin{tabularx}{\textwidth}{>{\bfseries}p{1.3cm}p{3cm}X>{\raggedright\arraybackslash}p{2cm}}
    \toprule
    ID & Thème & User Story & Statut \\
    \midrule
    7 & CV Pool & En tant que TA, je peux importer et gérer des CV dans un vivier centralisé. & Réalisé \\
    8 & Extraction IA & En tant que TA, je peux obtenir automatiquement les informations structurées d'un CV importé. & Réalisé \\
    9 & Recherche & En tant que TA, je peux rechercher des profils par texte, compétences et signaux sémantiques. & Réalisé \\
    10 & Détection de doublons & En tant que TA, je peux identifier des doublons potentiels dans le vivier. & Réalisé \\
    11 & Export documentaire & En tant que TA, je peux exporter un ou plusieurs CV depuis le vivier. & Réalisé \\
    12 & Gestion des offres & En tant que TA, je peux créer et visualiser des offres d'emploi structurées. & Réalisé \\
    13 & Templates d'offres & En tant que TA, je peux réutiliser une offre comme modèle. & Réalisé \\
    14 & Préparation du matching & En tant que TA, je peux relier les CV et les offres pour préparer le matching et l'affectation au pipeline. & Réalisé \\
    \bottomrule
  \end{tabularx}
\end{table}

La composition du backlog montre bien que ce sprint ne se limite pas à stocker des fichiers. Il couvre déjà la qualité documentaire, la recherche, l'export, la structuration des besoins et la préparation des futurs traitements intelligents.

\section{Implémentation du sprint}

\subsection{Analyse du sprint}
Le sprint s'appuie principalement sur les éléments fonctionnels suivants :
\begin{itemize}
  \item la page du vivier de CV ;
  \item l'interface de gestion et de consultation des profils ;
  \item la file d'upload en arrière-plan ;
  \item le service d'ingestion, d'extraction et de recherche des CV ;
  \item le service de détection des doublons ;
  \item le service de génération des embeddings ;
  \item le service de découpage des CV en chunks orientés RAG ;
  \item la page de gestion des offres ;
  \item l'interface de création et de listing des offres ;
  \item le service métier de gestion des offres ;
  \item la page de détail d'une offre et l'espace de pilotage associé au matching.
\end{itemize}

Ce module construit simultanément le patrimoine documentaire RH et la représentation structurée de la demande de recrutement. Il constitue donc la vraie base de travail des modules d'IA.

Il installe aussi un principe de robustesse important : l'import principal ne doit pas être bloqué par le temps de calcul des traitements sémantiques secondaires. Le système préfère livrer rapidement un résultat utile, puis compléter l'enrichissement en arrière-plan.
\paragraph{Qualité des données et transformation documentaire.}
L'un des enjeux majeurs du sprint concerne la qualité de la donnée d'entrée. Un CV importé peut être rédigé dans des styles très différents, contenir une mise en page complexe, mélanger plusieurs langues ou ne pas exposer les mêmes champs d'information qu'un autre document. Le système doit donc absorber une forte hétérogénéité entre PDF et DOCX, tout en acceptant qu'une extraction reste parfois partielle lorsque certaines informations sont absentes, ambiguës ou difficilement récupérables automatiquement.

Cette contrainte explique plusieurs choix du module : conservation du texte brut, extraction progressive, enrichissement différé, détection des doublons et séparation entre stockage principal et indexation sémantique. Dans la logique du rapport, ce point est essentiel, car la qualité du matching et du RAG dépend directement de la qualité documentaire initiale. Une base bruitée, redondante ou mal structurée affaiblit mécaniquement la pertinence des traitements IA qui s'appuient ensuite sur elle.

\paragraph{Tableau de transformation documentaire.}
Le passage du document brut à l'objet métier exploitable peut être résumé ainsi :

\begin{table}[H]
  \centering
  \caption{Transformation documentaire du CV et de l'offre}
  \begin{tabularx}{\textwidth}{>{\bfseries}p{3cm}X X}
    \toprule
    Entrée brute & Traitement & Sortie métier \\
    \midrule
    CV PDF ou DOCX hétérogène & Parsing du fichier et extraction du texte brut & Document lisible par le système \\
    Texte brut du CV & Extraction structurée des signaux métier : identité, compétences, expériences, formation, langues, résumé & Profil exploitable dans le vivier \\
    Profil extrait & Vérification de doublons, normalisation partielle et enrichissement progressif & Donnée documentaire plus fiable \\
    Profil enrichi & Génération d'un embedding global & Support de recherche sémantique directe \\
    Sections du profil & Découpage en chunks thématiques, embeddings par chunk et index lexical & Support du retrieval hybride et du futur RAG \\
    Offre structurée & Validation, normalisation des critères et stockage métier & Objet comparable avec les profils du vivier \\
    \bottomrule
  \end{tabularx}
\end{table}

Ce tableau montre que le CV n'est plus un simple fichier et que l'offre n'est plus un simple formulaire. Tous deux deviennent des objets métier progressivement transformés pour être réutilisables par les recruteurs, comparables par les services de matching et mobilisables par les couches RAG et agentiques.


\subsection{Vue fonctionnelle globale}
Les usages couverts par ce sprint se regroupent en trois blocs.
\begin{enumerate}
  \item \textbf{Gestion des CV} : importer un document, parser son contenu, extraire les informations principales, stocker le profil dans le vivier, le consulter, le rechercher, le supprimer et l'exporter.
  \item \textbf{Gestion des offres} : créer une offre, la décrire avec ses contraintes métier, la réutiliser comme modèle et la fermer de manière contrôlée.
  \item \textbf{Préparation de l'intelligence de recrutement} : générer des embeddings, découper les CV en chunks, rendre possible la recherche sémantique et le RAG, puis préparer la relation entre un CV, une offre et un futur candidat.
\end{enumerate}

Cette structuration montre que la plateforme traite à la fois l'offre et le profil comme deux objets complémentaires devant être rendus comparables par la machine sans perdre leur lisibilité métier.

\subsection{Pipeline d'upload et d'enrichissement d'un CV}
Le flux d'upload est l'un des plus importants du sprint. Il passe par \sourcecode{uploadCvAction()} puis par les services du module \sourcecode{cv-pool}. La logique est conçue pour que l'ingestion principale ne soit pas bloquée par les traitements avancés. La séquence peut être résumée ainsi :
\begin{itemize}
  \item le fichier est envoyé via la file d'upload côté client ;
  \item le service enregistre le CV dans \sourcecode{cv\_pool} ;
  \item le document est parsé selon son format ;
  \item l'IA extrait les informations structurées ;
  \item les champs métier sont mis à jour ;
  \item la génération de l'embedding est déclenchée en arrière-plan ;
  \item la génération des chunks pour le RAG est déclenchée en arrière-plan ;
  \item la vérification de doublons est exécutée ;
  \item l'interface se rafraîchit avec le résultat.
\end{itemize}

Cette architecture traduit un choix de robustesse : l'utilisateur n'attend pas que tout le pipeline IA soit terminé pour bénéficier de l'import principal. Elle traduit aussi un choix de qualité progressive : le CV devient d'abord exploitable, puis s'enrichit ensuite avec une profondeur sémantique plus fine.

\paragraph{Lecture opérationnelle du parcours.}
Un parcours représentatif illustre l'apport du module : un recruteur dépose un CV, le système extrait d'abord le texte et les métadonnées principales, puis le profil devient consultable dans le vivier. Les traitements plus lourds, comme l'embedding et le découpage en chunks, enrichissent ensuite le document pour les usages de recherche sémantique et de RAG. Cette séparation entre disponibilité rapide et enrichissement progressif évite de confondre ingestion documentaire et calcul IA complet.

\begin{figure}[H]
  \centering
  \diagraminput{figures/diagrams/cv-upload-pipeline}
  \caption{Séquence d'upload, d'extraction et d'indexation d'un CV}
\end{figure}

\subsection{Indexation sémantique et préparation du RAG}
Après extraction, le projet génère deux formes d'indexation : un embedding global du CV dans \sourcecode{cv\_pool} et un ensemble de chunks spécialisés dans \sourcecode{cv\_chunks}. L'embedding global sert à la recherche sémantique directe sur les profils. Les chunks, eux, servent à une granularité plus fine adaptée au RAG.

Le service \sourcecode{chunking.ts} découpe notamment les informations selon plusieurs sections : expérience, compétences, formation, résumé et langues. Chaque chunk reçoit ensuite son propre embedding et un index \sourcecode{tsvector}, ce qui prépare la récupération hybride combinant similarité vectorielle et recherche lexicale.

Cette séparation entre représentation globale et représentation fragmentée est centrale. Elle permet à la fois de retrouver un profil proche d'une intention de recherche, puis de justifier ce rapprochement à l'aide d'extraits précis et citables.

\paragraph{Différenciation des niveaux de recherche et d'indexation.}
Pour clarifier la logique du module, il est utile de distinguer cinq niveaux complémentaires :
\begin{itemize}
  \item la recherche simple, fondée sur des champs directement lisibles comme le nom, l'email, le fichier ou certaines compétences ;
  \item la recherche sémantique, qui s'appuie sur l'embedding global du CV pour rapprocher des profils au-delà des correspondances littérales ;
  \item le chunking, qui ne sert pas à rechercher mieux un profil complet, mais à découper son contenu en unités documentaires plus précises ;
  \item la préparation du RAG, qui exploite ces chunks, leurs embeddings et leur index lexical pour récupérer des passages pertinents avec citations ;
  \item le retrieval hybride, qui combine les signaux lexicaux et vectoriels afin de limiter les angles morts d'une recherche purement textuelle ou purement sémantique.
\end{itemize}

Autrement dit, l'embedding global répond surtout à la question : quel profil ressemble à cette intention de recherche ? Les chunks répondent plutôt à la question : quels passages précis du profil justifient cette réponse ? Cette distinction est structurante pour tout ce qui suivra dans le matching, le RAG et l'agent conversationnel.

\subsection{Création d'une offre et affectation vers le pipeline candidat}
La création d'une offre passe par \sourcecode{createJobAction()} puis \sourcecode{createJob()}. Les données sont validées avant insertion. L'offre conserve une structure explicite : titre, description, compétences indispensables, compétences souhaitables, séniorité, business unit, statut et indicateur de template.

Même si le cœur du pipeline candidat est traité dans le chapitre suivant, ce module prépare déjà le lien métier principal : le passage de \textit{CV dans le vivier} à \textit{candidat dans une offre}. Cette transition transforme un actif documentaire en objet de pipeline.

Elle rend aussi possible un futur matching crédible, car la plateforme dispose désormais de deux objets structurés pouvant être comparés selon des critères métier et sémantiques.

\begin{figure}[H]
  \centering
  \diagraminput{figures/diagrams/job-assignment-sequence}
  \caption{Création d'un poste et affectation d'un CV au pipeline candidat}
\end{figure}

\section{Réalisation}

\subsection{Page \textit{CV Pool}}
La page du vivier de CV charge la liste des CV ainsi que les statistiques du vivier via \sourcecode{listCvPoolAction()} et \sourcecode{getCvPoolStatsAction()}. Le composant \sourcecode{CvPoolClient} fournit une interface riche qui permet :
\begin{itemize}
  \item l'upload par sélection ou glisser-déposer ;
  \item l'usage d'une file globale d'upload ;
  \item la recherche locale sur nom, email, fichier et compétences ;
  \item la sélection multiple ;
  \item la suppression unitaire ou par lot ;
  \item l'ouverture d'un dialogue de revue détaillée du CV extrait ;
  \item le téléchargement du fichier brut ;
  \item l'export individuel Excel et l'export multiple ZIP.
\end{itemize}

La page présente aussi un niveau de synthèse à travers \sourcecode{CvPoolStats} : nombre total de CV, top compétences, distribution des langues et tendance d'upload sur les sept derniers jours.
La version actuelle va plus loin qu'une simple revue documentaire : le dialogue de consultation détaillée embarque désormais un bloc d'\textit{evidence readiness} qui explicite les faits observés, les données manquantes et les risques documentaires avant toute action IA. Les actions \textit{Explain CV with Agent} et \textit{Draft Questions with Agent} ouvrent ensuite le workspace \sourcecode{/agent} avec un prompt prérempli qui réinjecte ce contexte de preuve plutôt que de repartir d'une page blanche.
Cette interface montre ainsi que le vivier n'est pas une simple liste de fichiers. C'est une base de profils analysés, consultables et réexploitables, avec une lecture à la fois opérationnelle, analytique et déjà branchée sur l'assistance agentique.

\begin{figure}[H]
  \centering
  \screenshotplaceholder{CV pool}
  \caption{Page de gestion du vivier de CV.}
\end{figure}

\begin{figure}[H]
  \centering
  \screenshotplaceholder{CV detail dialog}
  \caption{Dialogue de consultation détaillée d'un CV.}
\end{figure}

\subsection{Upload en arrière-plan et robustesse d'ingestion}
Le composant \sourcecode{upload-provider.tsx} introduit une file d'upload avec états \sourcecode{queued}, \sourcecode{uploading}, \sourcecode{success} et \sourcecode{error}, validation de format, limite de taille, \textit{retries} et \textit{backoff} exponentiel. Cette logique améliore nettement l'expérience utilisateur et la robustesse du module, car elle absorbe les traitements asynchrones sans bloquer l'usage.

Elle répond à une contrainte très concrète : un recruteur peut déposer plusieurs CV à la suite et continuer à travailler pendant que l'enrichissement documentaire progresse en tâche de fond.

\subsection{Détection de doublons}
Le projet intègre un mécanisme de détection des doublons dans \sourcecode{duplicate-detection.ts}. Le service compare notamment l'email extrait, la similarité de nom et le numéro de téléphone. Cette fonctionnalité améliore la qualité documentaire du vivier et, par conséquent, la qualité du matching et de la recherche RAG qui s'appuient sur ces données.

La qualité de la base est ici un prérequis de la qualité algorithmique. Une couche IA branchée sur des profils redondants ou incohérents produirait mécaniquement des résultats moins fiables.

\subsection{Page \textit{Jobs}}
La page de gestion des offres s'appuie sur \sourcecode{listJobsAction()} et \sourcecode{getJobsStatsAction()}. Le composant \sourcecode{JobsClient} combine des cartes statistiques, la liste des offres, un dialogue de création d'offre, la visualisation des statuts et des signaux sur la demande en compétences et en séniorité.

Les statistiques associées incluent le total des jobs, la répartition par séniorité, par statut, par business unit et les compétences les plus demandées. Cette page transforme les besoins recruteurs en objets comparables, mesurables et exploitables par les services de matching.

Elle joue aussi un rôle de normalisation : en imposant une structure cohérente à la formulation des postes, elle facilite les analyses globales et la réutilisation ultérieure dans les traitements IA.

L'offre d'emploi n'est donc pas traitée comme un simple formulaire administratif. Elle devient un objet de décision qui articule besoin métier, contraintes de recrutement et future lecture algorithmique. C'est précisément cette structuration qui permet ensuite de comparer de manière cohérente un poste, un profil, un screening et un résultat de matching.

\begin{figure}[H]
  \centering
  \screenshotplaceholder{Job list}
  \caption{Catalogue des offres d'emploi.}
\end{figure}

\begin{figure}[H]
  \centering
  \screenshotplaceholder{Create/edit job area}
  \caption{Zone de création ou d'édition d'une offre d'emploi.}
\end{figure}

\subsection{Templates et fermeture contrôlée des offres}
Le service \sourcecode{jobs.ts} supporte l'enregistrement d'une offre comme template, la liste des templates, la création d'une offre à partir d'un template et la fermeture contrôlée d'un poste. Cette fermeture est sécurisée par des règles métier : elle est interdite s'il reste des entretiens planifiés ou des candidats encore actifs dans le pipeline. Cette règle protège la cohérence fonctionnelle du système.

Elle illustre une approche importante du projet : une action d'administration simple en apparence reste soumise à l'état réel du workflow métier.

\subsection{Page détail d'une offre}
La page de détail d'une offre charge l'offre, les candidats qui lui sont liés et la liste des managers. Le composant \sourcecode{JobDetailClient} transforme ensuite cette page en espace de pilotage avancé, avec génération de screening IA, consultation du score, génération et édition des questions d'entretien, planification, envoi d'emails, mise à jour des étapes candidat, assignation à un manager et matching \textit{inline} de CV.
L'interface distingue désormais explicitement les actions directes des actions de raisonnement IA. L'ouverture du matching de CV et celle des entretiens reposent sur des onglets pilotés par l'URL, ce qui rend les liens profonds robustes et permet au recruteur d'atterrir directement sur le bon sous-espace. À côté de cela, les entrées agentiques restent dédiées à l'explication du matching, à la revue d'exigences et à la synthèse raisonnée.
Ce point montre que ce module n'est pas isolé. Le détail d'une offre devient le point de convergence entre la demande de recrutement, la base documentaire des profils, les parcours directs du produit et les futurs mécanismes de décision.

\begin{figure}[H]
  \centering
  \screenshotplaceholder{Job detail + matching}
  \caption{Page détail d'une offre avec candidats associés et résultats de matching.}
\end{figure}

\section{Apport du module}
Ce module apporte une contribution majeure au projet :
\begin{itemize}
  \item il transforme des documents bruts en objets structurés ;
  \item il crée la matière première nécessaire au matching et au RAG ;
  \item il structure la demande métier à travers des offres exploitables ;
  \item il améliore la qualité documentaire via la détection de doublons ;
  \item il prépare la relation fondamentale \textit{CV $\rightarrow$ offre $\rightarrow$ candidat} ;
  \item il installe les bases techniques de la recherche sémantique et de la récupération de contexte.
\end{itemize}

Dans la logique globale du rapport, ce module représente le moment où la plateforme commence réellement à devenir intelligente, car les données qu'elle manipule ne sont plus seulement visibles à l'écran : elles deviennent analysables, indexables et interrogeables par l'IA.

Il faut surtout retenir que cette intelligence n'apparaît pas magiquement dans le chapitre suivant. Elle devient possible ici, parce que la plateforme commence à transformer des fichiers bruts en structures de données propres, reliées et réexploitables.

\section*{Conclusion}
\addcontentsline{toc}{section}{Conclusion}
Ce deuxième module a permis de construire le vivier documentaire et la base métier du recrutement. Les CV sont désormais importés, analysés, enrichis, exportables et progressivement indexés pour la recherche sémantique et le RAG. Les offres d'emploi sont quant à elles structurées de manière cohérente pour soutenir le matching et les décisions de recrutement. Le chapitre suivant s'appuiera sur cette base pour traiter la progression des candidats dans le pipeline et la coordination des entretiens.
